\scanpage{18}

\mspage{\handwritten{3}}

``wholly degen\deleted{\unsure{r}}erate'' generalized boundary). When we throw in the rigidifications and go over to isomorphism classes, this should give the desired retraction. However, the geometric situation is a lot richer still, as the compact surface without boundary obtained through surgery is endowed, not only with a system of \handwritten{$\nu$} points, b\deleted{\unsure{y}}ut moreover with a system of mutually disjoint \uline{discs} around these points. The shape of these discs is by no means arbitrary - we'll say that a system of discs around \handwritten{$\nu$} points on a compact \inserted{conformal} surface $X^{\wedge}$ without boundary is ``admissible'', if the situation can be obtained as above (up to isomorphism) from surgery, starting with a compact conformal surface $X$ w\deleted{\unsure{j}}ith bo\uline{u}ndary. (NB Among the given``discs'', we should allow \deleted{\unsure{s}}that some should be reduced to their center - we'll call them ``degenerate''\handwritten{)}The condition of admissibility can be expressed intrinsically, by stating that \deleted{the} for every non-degenerate component $\Gamma_{\mathrm i}$ of the system of boundaries of those discs, the two operations we got of the standard circle group (of complex numbers of module 1) upon $\Gamma_{\mathrm i}$, by using the fact that it is (on the one hand) the boundary of the disc $D_{\mathrm i}$, and (on the other hand) that it is a component of the boundary of the hyperbolic surface $X^\wedge\backslash(\handwritten{\bigcup_{\mathrm j}}D^{\mathrm o}_{\mathrm j})$, should be the same. \margin{\handwritten{!}} When $X^\wedge$ and the points $s_{\mathrm i}$ on $X^\wedge$ are given, the possible \inserted{admissible} systems of discs around the points $s_{\mathrm i}$ depend on \handwritten{$\nu$} parameters - and the first idea which flips to mind to give a more precise meaning to these ``parameters'', is to view them as being the``radii'' of those discs. But then we'll have to define what we mean by these !

The idea here is that, when we have a conformal disc $D$ and an inter\deleted{\unsure{a}}ior point $s$ of $D$, then $D$ may be viewed as canonically embedded in \deleted{\unsure{h}}the tangent space $T_{\mathrm s}$ to $D$ at $s$, as th\^{e} ``unit disc'' at $s$. Thus, in the situation above of an\deleted{\unsure{just}} admissible system of discs $(D_{\mathrm i})_{\mathrm i\handwritten{\in}\mathrm I}$ around $(s_{\mathrm i})_{\mathrm i\handwritten{\in}\mathrm I}$, for every $s_{\mathrm i}$ corresponding to a non-degenerate $D_{\mathrm i}$, we get a canonical disc
\[
\handwritten{\Delta}_{\mathrm i}\handwritten{\subset} T_{s_{\mathrm i}}
\]
in the tangent space - and of course, for degenerate $D_{\mathrm i}$, we'll take $\handwritten{\Delta}_{\mathrm i}$ to be degenerate too. The discs we get in a\deleted{g} given $T_{s_{\mathrm i}}$ (for a fixed system $(s_{\mathrm i})$, and a variable \inserted{admissible} system of discs around these $s_{\mathrm i}$) are all discs in the strict euclidean sense, given by an unequality
\[
\handwritten{|}z\handwritten{|}\ \handwritten{\leqslant}\ r_{\mathrm i}\quad,
\]
where $z\handwritten{\mapsto|z|}$ denotes some hermitian metric on $T_{s_{\mathrm i}}$ compatible with the conformal structure - this metric being \deleted{defined}\inserted{\handwritten{unique}} up to a scalar factor. the set \inserted{$\mathrm R_{\mathrm i}$} of all those possible discs (the non-degenerate one\handwritten{s}\deleted{\unsure{,}}, say) may be viewed in a natural way as a ``torsor'' (= principal homogeneous space)
% archive sheet number: 17
